\exnumber{01}
\extitle{UI III - Ex 2 - Multi state}
\exfiles{ex01/Figma file, named Name\_FirstName\_DeliverableName\_Date\_VersionNumber}
\exauthorize{Those named in the exercise}
\makeheaderfiles

\textbf{Instructions}:

Go to Sony.com \url{http://sony.com} and observe the buttons and call to actions. When you hover or click on them, they change!

Make screenshots of the different states of buttons, menu titles, filters etc. Then, gather and organize them.

Find at least five combinations of multi state elements.

\textbf{Deliverables:}

\begin{itemize}
\item A Figma file with 5 combinations of multi state elements.
\item Upload the Figma file on Github, labeled as follows: "Name\_FirstName\_DeliverableName\_Date\_VersionNumber". For example, it can be labeled as "Doe\_John\_MidFiPrototype\_05242024\_V1".
\item Add a title within the file.
\item Pay attention to the global aspect of the file. It must be clear, simple and easy to read. You can get inspiration on canva.com but don't overload your design with too much details!
\end{itemize}

\textbf{Resources:}

\begin{itemize}
\item Video tutorial on Figma \url{https://www.youtube.com/watch?v=FTFaQWZBqQ8}
\item Beginning graphic design \url{https://edu.gcfglobal.org/en/beginning-graphic-design/}
\item Gareth graphic studio \url{https://www.youtube.com/c/GarethDavidStudio}
\item Example here with Nintendo.com \url{https://www.figma.com/file/ahSnWRZeKBO62oJDiXltxY/UI-III---Ex-2}
\end{itemize}

/!\textbackslash{} The screenshots date back to nov 2021. They may differ from the current version of Nintendo.com \url{http://Nintendo.com} but the rules are the same

